\documentclass[12pt,reqno]{amsart}
\usepackage{amsmath,amssymb}
\usepackage{mathptmx}
\usepackage[T1]{fontenc}
\usepackage{enumitem}
\pagestyle{plain}

\setlength{\oddsidemargin}{0in}
\setlength{\evensidemargin}{0in}
\setlength{\topmargin}{0in}
\setlength{\headheight}{0pt}
\setlength{\headsep}{0pt}
\setlength{\footskip}{.3in}
\setlength{\textheight}{9.2in}
\setlength{\textwidth}{6.5in}
\setlength{\parindent}{0pt}
\setlist[enumerate,1]{leftmargin=*,itemsep=.9em}
\setlist[enumerate,2]{label=(\alph*),leftmargin=1.7em,itemsep=.3em,topsep=.25em}

\newcommand{\C}{\mathbb{C}}

\title{Problem Set 5 \\ MATH 411: Honors Complex Variables, Fall 2026}

% New exercises patterned after Fall 2021 Homeworks 5 and 6.
% Coverage: Lectures 23--25.

\begin{document}
\maketitle

These are practice problems and are not to be turned in. Justify your
answers and all limiting arguments. Unless stated otherwise, closed
contours are oriented counterclockwise.

\begin{enumerate}

\item \textbf{Contour integrals by residues.}
Locate and classify the isolated singularities of the integrand, then
evaluate
\[
 \int_{|z|=r}\frac{e^z-1}{z^3(z-2)}\,dz
 \qquad\text{for }r=1\text{ and }r=3.
\]

\item \textbf{A trigonometric integral.}
Use the substitution $z=e^{it}$ and the residue theorem to evaluate
\[
 \int_0^{2\pi}\frac{\cos(2t)}{5-4\cos t}\,dt.
\]
\emph{Hint.} Write $\cos(2t)=\operatorname{Re}(e^{2it})$.

\item \textbf{An improper integral.}
For $a>0$, use an upper-half-plane contour to evaluate
\[
 \int_{-\infty}^{\infty}\frac{\cos(ax)}{(x^2+4)^2}\,dx.
\]
Give an explicit bound for the integral of $e^{iaz}/(z^2+4)^2$
over the semicircular arc $|z|=R$, $\operatorname{Im}z\geq0$, for $R>2$.

\item \textbf{The argument principle.}
For each $r>0$, let $C_r$ be the circle $|z|=r$. Consider
\[
 F(z)=\frac{(e^z-1)^2}{z^3(z-1)^2(z-3)}.
\]
\begin{enumerate}
\item Find all zeros and poles of $F$ in $|z|<7$, with their orders.
\item Use the argument principle to find the winding numbers $n_2$ and
$n_7$ of the image curves of $C_2$ and $C_7$ under $F$ about $0$.
What does $n_7-n_2$ count in the annulus $2<|z|<7$?
\end{enumerate}

\item \textbf{Rouch\'e's theorem.}
\begin{enumerate}
\item Let $a\in\C$ with $|a|\leq2$. Show that
$P_a(z)=z^6-6z^2+az+1$ has no zeros on $|z|=1$ or $|z|=2$.
How many zeros, counted with multiplicity, lie in each of the regions
$|z|<1$, $1<|z|<2$, and $|z|>2$? Justify the counts using
Rouch\'e's theorem.
\item Use Rouch\'e's theorem to prove the Fundamental Theorem of Algebra:
every complex polynomial of degree $n\geq1$ has exactly $n$ zeros in
$\C$, counted with multiplicity. \emph{Hint.} Compare the polynomial
with its leading term on a sufficiently large circle.
\end{enumerate}

\item \textbf{The open mapping theorem.}
Let $U\subseteq\C$ be a bounded domain and let $f$ be continuous on
$\overline U$ and holomorphic and nonconstant on $U$.
\begin{enumerate}
\item Use the open mapping theorem and compactness to prove that
$\partial(f(U))\subseteq f(\partial U)$.
\item For $U=\{z:|z-1|<2\}$ and $f(z)=z^2$, show that the inclusion is strict.
\end{enumerate}

\end{enumerate}
\end{document}
